% Independent audit, 2026-10-09. Notes for integration by the parent agent.
\section{Exact rational certificates for Gaussian values}\label{sec:certificates}

We use the outermost-letter-first convention
\[
 G(a_1,\ldots,a_w;z)=\int_0^z
       \frac{G(a_2,\ldots,a_w;t)}{t-a_1}\,dt,
 \qquad G(;z)=1.
\]
Define the permitted alphabet region
\[
 \mathcal A=\{0,1\}\cup
 \{a\in\mathbb Q(i): |a|\geq1,\ |1-a|\geq1\}.
\]
The two inequalities are exact rational comparisons after squaring the
moduli. The Gaussian alphabet and the letter $1-i$ both belong to this
region.

\begin{theorem}[Uniform exact certificate]\label{thm:certificate}
Let $w\geq1$, let $a_1,\ldots,a_w\in\mathcal A$, and assume
$a_1\ne1$ and $a_w\ne0$. Then $G(a_1,\ldots,a_w;1)$ exists as an ordinary
convergent iterated integral. It has the Hölder decomposition
\begin{equation}\label{audit:eq:holder}
 G(a_1,\ldots,a_w;1)
 =\sum_{j=0}^w(-1)^j
  G(1-a_j,\ldots,1-a_1;1/2)
  G(a_{j+1},\ldots,a_w;1/2),
\end{equation}
where a missing word denotes $1$.
Replace each factor by its Taylor polynomial through degree $N$, evaluated
at $1/2$, and call the resulting Gaussian rational number $C_N$. Then
\begin{equation}\label{audit:eq:radius}
 |G(a_1,\ldots,a_w;1)-C_N|
 \leq 2(w+1)\,T_w(N),\qquad
 T_d(N)=2^{-N}\sum_{k=0}^{d-1}\binom Nk.
\end{equation}
Here $\binom Nk=0$ for $k>N$ and $T_0(N)=0$.
All the Taylor polynomials required by one word are computed using
$O(wN)$ exact rational arithmetic operations.
For an integer $P\geq0$, the explicit choice
\begin{equation}\label{audit:eq:Nchoice}
 N=3\bigl(P+w+\lceil\log_2(w+1)\rceil\bigr)
\end{equation}
guarantees an error at most $2^{-P}$. In particular, $N=O(P+w)$.
\end{theorem}

\begin{proof}
The hypotheses exclude singularities in the open integration segment.
They also ensure integrability at its endpoints: the last letter is
nonzero, and any logarithmic singularities at $1$ are integrated against
an outermost kernel regular there. Partition the ordered integration
simplex according to the number $j$ of variables above $1/2$. Replacing
those variables by their complements reverses their word and contributes
$(-1)^j$. This gives \eqref{audit:eq:holder}; this is the standard Hölder
convolution, with its endpoint hypotheses made explicit.

Every nonzero letter in either factor of \eqref{audit:eq:holder} has
modulus at least $1$. Moreover, every nonempty right factor ends at
$a_w\ne0$, and every nonempty left factor ends at $1-a_1\ne0$.
Thus these factors have Taylor series with zero constant coefficient.

Consider a word $v$ of this kind, with $d$ nonzero letters, and write
$G(v;z)=\sum_{n\geq0}c_nz^n$. Coefficientwise absolute values are bounded
by replacing each nonzero kernel with $1/(1-z)$. Inserting a zero kernel
divides the coefficient of $z^n$ by $n\geq1$, so deleting all zero letters
can only increase this majorant. Consequently
\[
 |c_n|\leq[z^n]\frac{(-\log(1-z))^d}{d!}
 \leq\binom{n-1}{d-1}\quad(n\geq d),
 \qquad c_n=0\quad(n<d).
\]
For the second inequality, expand the $d$ logarithms; each composition of
$n$ into $d$ positive parts contributes at most $1$, and the additional
factor $1/d!$ can be discarded. The empty word is handled separately.
It follows that
\[
 |G(v;1/2)|\leq1,\qquad
 \left|\sum_{n=0}^Nc_n2^{-n}\right|\leq1,
\]
because $\sum_{n\geq d}\binom{n-1}{d-1}2^{-n}=1$.
The tail of this majorant is exactly
\[
 \sum_{n>N}\binom{n-1}{d-1}2^{-n}
 =2^{-N}\sum_{k=0}^{d-1}\binom Nk=T_d(N).
\]
One proof interprets the left side as the probability that the $d$th
success in fair independent Bernoulli trials occurs after trial $N$;
the right side counts fewer than $d$ successes among the first $N$
trials. Equivalently, the equality follows by Pascal's identity.

For each product $AB$ in the Hölder sum, with truncations $A_N,B_N$,
\[
 |AB-A_NB_N|
 \leq |A-A_N|\,|B|+|A_N|\,|B-B_N|
 \leq T_{d_A}(N)+T_{d_B}(N)\leq2T_w(N).
\]
Summing over the $w+1$ cuts proves \eqref{audit:eq:radius}.

For the complexity claim, let $b_n$ be the Taylor coefficients of a
tail and $c_n$ those of the word obtained by prepending $a$. For $n\geq1$,
\begin{equation}\label{audit:eq:recurrence}
 c_n=
 \begin{cases}
   b_n/n,&a=0,\\
   \bigl((n-1)c_{n-1}-b_{n-1}\bigr)/(an),&a\ne0.
 \end{cases}
\end{equation}
In both cases $c_0=0$, and the empty tail has $b_0=1$, $b_n=0$ for
$n\geq1$. The second recurrence is coefficient comparison in
$(z-a)G'(a,v;z)=G(v;z)$.
It takes $O(N)$ operations per prepended letter. Compute the original
suffix chain starting from the last letter; this gives every right
factor. Compute the reflected reversed-prefix chain by successively
prepending $1-a_1,1-a_2,\ldots,1-a_w$; this gives every left factor.
Each chain requires $O(wN)$ operations in total. Summing the $w+1$
products adds $O(w)$ operations.

Finally, the binomial theorem gives the useful bound
\[
 T_w(N)\leq 2^{w-1}(3/4)^N.
\]
Indeed, for $k\leq w-1$ one has $2^{-k}\geq2^{-(w-1)}$, so
\[
 \sum_{k=0}^{w-1}\binom Nk
 \leq2^{w-1}\sum_{k=0}^N\binom Nk2^{-k}
 =2^{w-1}(3/2)^N.
\]
With $M=P+w+\lceil\log_2(w+1)\rceil$ and $N=3M$,
\[
 2(w+1)T_w(N)
 \leq(w+1)2^w(3/4)^{3M}
 \leq(w+1)2^{w-M}\leq2^{-P},
\]
since $27/64<1/2$. This proves \eqref{audit:eq:Nchoice}.
\end{proof}

\begin{corollary}[Fixed imaginary quadratic fields]
\label{cor:quadratic-certificate}
Let $K\subset\mathbb C$ be an imaginary quadratic field with its fixed
complex embedding. Replace $\mathbb Q(i)$ in the alphabet definition by
$K$. The same H\"older sum, tail bound, precision prescription, and
$O(wN)$ arithmetic-operation count hold, with the center now in $K$.
\end{corollary}
\begin{proof}
The proof of Theorem~\ref{thm:certificate} uses the coefficient field
only for exact additions, multiplications, and reciprocals. Each
Taylor coefficient therefore lies in $K$, and the analytic majorant
is unchanged. The tests $|a|^2\ge1$ and $|1-a|^2\ge1$ are exact
rational comparisons because these squared moduli are field norms.
Arithmetic in a fixed quadratic field requires a bounded number of
rational operations per field operation.
\end{proof}

In particular, every sixth root of unity satisfies the alphabet
condition: apart from $1$, its distance from $1$ is at least $1$.
Thus the theorem also supplies an exact-certificate route for the
Eisenstein values. The supplied core implementation specializes to
$\mathbb Q(i)$; the four mixed Eisenstein formulas are established
analytically and checked numerically in the independent scripts.

\paragraph{A safe bit-complexity statement.}
The count $O(wN)$ is an arithmetic-operation count; it does not assign
unit cost to arbitrarily large integers. Let $Q$ be a common positive
denominator for the real and imaginary components of the reciprocals
of all nonzero letters used in both Hölder chains. Let
$L_N=\operatorname{lcm}(1,\ldots,N)$. For a word of length $\ell\leq w$,
each component of its degree-$n$ Taylor coefficient has denominator
dividing $Q^nL_n^\ell$. This follows from the convolution form of
\eqref{audit:eq:recurrence}: a term uses total reciprocal-letter degree
$n$ and one integer divisor at most $n$ for each integration.
Thus all evaluated partial sums have denominators dividing
\[
 D_N=(2Q)^N L_N^w.
\]
The coefficient majorant bounds their numerator bit lengths by
\[
 B=O\bigl(N\log(2Q)+wN\log(N+1)+N\bigr),
\]
using only the elementary estimate $L_N\leq N!$.
All product sums have denominator dividing $D_N^2$ and comparable
operand lengths. Therefore standard exact rational arithmetic yields
polynomial bit complexity in $w,N,\log Q$. For an entirely conservative
explicit estimate, naive multiplication and division use $O(B^2)$ bit
operations, the Euclidean algorithm uses at most $O(B)$ such divisions,
and normalization after every operation yields $O(wNB^3)$ bit
operations. This deliberately loose bound requires no fast-arithmetic
assumptions. Standard fast multiplication and gcd algorithms improve it.

\paragraph{What is and is not new here.}
Hölder convolution and its use for numerical multiple-polylogarithm
evaluation are established; see Borwein--Bradley--Broadhurst--Lison\v ek~\cite{BBBL2001}
and Vollinga--Weinzierl~\cite{VollingaWeinzierl2005}. The contribution of this specialization is the
explicit rational tail certificate, the permitted rational Gaussian
letter region, the reuse of both chains, and the uniform linear bound
$N=O(P+w)$. This is an error bound for computing periods, not a proof of
an identity inferred from a small residual.

\paragraph{Application to the manuscript's constants.}
The letters needed for $\pi$, $\log2$, $\zeta(n)$ for $n\geq2$,
$\beta(n)$, and $\lambda_n=\Im\Li_n((1+i)/2)$ all satisfy the hypotheses:
\[
 \pi=4\Im\Li_1(i),\quad \log2=-2\Re\Li_1(i),\quad
 \Li_n(z)=-G(0^{n-1},z^{-1};1).
\]
For the half point the nonzero letter is $1-i$, whose reflected letter
is $i$. Hence both sides of the Gaussian conjectures can be enclosed
using one exact Gaussian-rational arithmetic framework.

\subsection{Agreement with ordinary boundary series}

For completeness, suppose $s_1,\ldots,s_d$ are positive integers and
$|z_j|=1$. The outer-first colored multiple-polylogarithm series converges
absolutely if $s_1\geq2$, and converges ordinarily if $s_1=1$ and
$z_1\ne1$. These sufficient conditions include every colored value used
in the Gaussian tables.

For $d\geq2$, put
\[
 B_n=\sum_{n>n_2>\cdots>n_d\geq1}
       \prod_{j=2}^d\frac{z_j^{n_j}}{n_j^{s_j}}.
\]
The elementary symmetric harmonic bound gives
$|B_n|=O((1+\log n)^{d-1})$ and
\[
 |B_{n+1}-B_n|
   =O\bigl((1+\log n)^{d-2}/n\bigr).
\]
The first estimate proves absolute convergence for $s_1\geq2$.
For $s_1=1$ set $b_n=B_n/n$. Then $b_n\to0$ and
\[
 |b_{n+1}-b_n|
    =O\bigl((1+\log n)^{d-1}/n^2\bigr),
\]
which is summable. The geometric partial sums of $z_1^n$ are bounded
when $z_1\ne1$, so summation by parts proves convergence of
$\sum z_1^n b_n$. At $d=1$ this is the usual Dirichlet test.

Introduce $0<r<1$ in the outer color, $z_1\mapsto rz_1$. The absolutely
convergent series/integral dictionary gives the GPL with the original
word and endpoint $r$. Abel's theorem identifies the limit of the
series with its ordinary boundary sum. On the integral side, the last
letter is nonzero. Near zero every nonempty suffix is analytic and
$O(t)$. Near one, recursive integration gives at worst powers of
$|\log(1-t)|$ in an interior suffix. The outer letter is either $0$
(when $s_1\geq2$) or $z_1^{-1}\ne1$ (when $s_1=1$), so its kernel is
bounded near one and these logarithmic powers are integrable. Thus the
GPL has a continuous limit as $r\uparrow1$ and agrees with the ordinary
boundary series. No regularization is being silently used in these
cases.
